\documentclass[UTF8,fontset=none,aspectratio=169,10pt]{beamer}
\usepackage{ads-beamer-theme}
\title{高级分布式系统}
\subtitle{第十六讲：分布式认证、安全与可信身份}
\author{陈鹏飞\quad 计算机学院}
\institute{中山大学}
\date{2025--2026 学年}

\begin{document}
\begin{frame}[plain]
\begin{tikzpicture}[remember picture,overlay]
  \fill[ADSBlue] (current page.north west) rectangle ([xshift=2.4cm]current page.south west);
  \fill[ADSOrange] ([xshift=2.4cm]current page.north west) rectangle ([xshift=2.55cm]current page.south west);
  \node[anchor=west,align=left] at ([xshift=3.3cm,yshift=1.0cm]current page.west) {\fontsize{32}{38}\selectfont\color{ADSBlue}\bfseries 高级分布式系统};
  \node[anchor=west,align=left] at ([xshift=3.3cm,yshift=-0.4cm]current page.west) {\fontsize{21}{27}\selectfont\color{ADSDarkBlue}第十六讲\quad 分布式认证与安全};
  \node[anchor=west,align=left] at ([xshift=3.3cm,yshift=-2.2cm]current page.west) {\large\color{ADSGray}系统如何证明“你是谁”，又如何限制“你能做什么”？};
  \node[anchor=east] at ([xshift=-0.8cm,yshift=-3.2cm]current page.east) {\small\color{ADSGray}陈鹏飞\quad 中山大学\quad 2025--2026};
  \draw[ADSBlue,line width=0.8pt] ([xshift=3.3cm,yshift=-1.35cm]current page.west) -- ([xshift=12.6cm,yshift=-1.35cm]current page.west);
\end{tikzpicture}
\end{frame}

\sectionframe{一、认证问题从哪里来？}{网络把“调用者、服务和消息”都暴露在不可信环境中。}

\begin{frame}{本讲学习目标}
\begin{enumerate}
  \item 区分 authentication、authorization、confidentiality、integrity；
  \item 理解共享密钥、公开密钥、nonce、时间戳和票据；
  \item 能手工分析重放攻击，并修复基本认证协议；
  \item 读懂认证服务器与 Kerberos 的三段式流程；
  \item 将协议原则映射到 mTLS、服务网格和 AI 基础设施。
\end{enumerate}
\vfill
\callout{主线}{先证明身份，再决定权限；协议必须同时回答“新鲜性、机密性、完整性”三个问题。}
\end{frame}

\begin{frame}{分布式系统中的 principal}
\begin{center}
\begin{tikzpicture}[node distance=12mm and 16mm,>=Latex]
  \nodebox[minimum width=2.3cm]{ADSBlue}{u}{用户/客户端};
  \nodebox[right=of u,minimum width=2.3cm]{ADSOrange}{svc}{服务};
  \nodebox[below=of u,minimum width=2.3cm]{ADSGreen}{idp}{身份服务};
  \nodebox[below=of svc,minimum width=2.3cm]{ADSGray}{att}{攻击者};
  \draw[\arrow,ADSBlue] (u)--node[above,font=\scriptsize]{请求}(svc);
  \draw[\arrow,ADSGreen] (idp)--node[left,font=\scriptsize]{签发/验证}(u);
  \draw[\arrow,ADSGreen] (idp)--node[right,font=\scriptsize]{策略}(svc);
  \draw[\arrow,ADSOrange,dashed] (att)--node[right,font=\scriptsize]{窃听/篡改/重放}(svc);
\end{tikzpicture}
\end{center}
\begin{itemize}
  \item principal 可以是用户、进程、服务账号、设备或控制器；
  \item “身份”是名字，\key{凭证}是证明名字的证据，\key{权限}是允许执行的动作。
\end{itemize}
\end{frame}

\begin{frame}{四个经常混淆的安全目标}
\begin{tabularx}{\linewidth}{>{\bfseries}p{3.0cm}X}
\toprule
认证 Authentication & 对方是谁？\quad 例如确认请求确实来自 scheduler-7。\\
授权 Authorization & 这个身份能做什么？\quad 例如只能读取模型版本，不得删除生产版本。\\
机密性 Confidentiality & 旁观者看不到内容。\\
完整性 Integrity & 消息在传输和存储过程中未被篡改。\\
\bottomrule
\end{tabularx}
\vfill
\callout{例子}{签名能证明“内容由谁签过”，但不自动说明“这个身份是否有权执行操作”。}
\end{frame}

\begin{frame}{威胁模型：攻击者能做什么？}
\begin{itemize}
  \item 窃听：记录用户名、票据、密文和时间；
  \item 篡改：替换参数、重排消息或修改路由；
  \item 冒充：声称自己是 P、Q 或某个服务账号；
  \item 重放：把过去的合法消息在未来再次发送；
  \item 拒绝服务：让身份服务不可用，或制造大量无效请求。
\end{itemize}
\vfill
\smallnote{本讲以“攻击者可完全控制网络，但不能直接读出正确保管的密钥”为基本假设。}
\end{frame}

\begin{frame}{单向认证与双向认证}
\begin{columns}[T,onlytextwidth]
\column{.48\linewidth}
\blocktitle{单向认证}
\begin{itemize}
  \item 客户端验证服务器；
  \item 典型：浏览器验证网站证书；
  \item 适合公开服务，但客户端身份可能仍未知。
\end{itemize}
\column{.48\linewidth}
\blocktitle{双向认证}
\begin{itemize}
  \item 客户端与服务器互相证明身份；
  \item 典型：mTLS、服务到服务调用；
  \item 多一步握手，却减少“假服务”和横向移动。
\end{itemize}
\end{columns}
\vfill
\callout{思考}{为什么浏览器通常不要求用户证书？因为网站需要识别的是“账户”，而不是直接识别每台终端。}
\end{frame}

\sectionframe{二、密码学工具箱}{不同工具解决不同问题：不要用“加密”这个词覆盖全部语义。}

\begin{frame}{对称密码与非对称密码}
\begin{tabularx}{\linewidth}{>{\bfseries}p{3.1cm}XX}
\toprule
对称密码 & P、Q 共享同一个秘密密钥；计算快；密钥分发困难。 & 批量数据加密、会话加密 \\
非对称密码 & 公钥公开，私钥保密；可签名和验签；计算成本较高。 & 身份证明、密钥交换、证书 \\
混合方案 & 用非对称密码建立会话密钥，再用对称密码传数据。 & TLS、Kerberos 的工程化变体 \\
\bottomrule
\end{tabularx}
\end{frame}

\begin{frame}{Hash、MAC、数字签名}
\begin{itemize}
  \item Hash$(x)$：固定长度摘要，适合检测变化，但没有身份含义；
  \item MAC$_k(x)$：只有持有共享密钥的人才能生成和验证，提供完整性与来源认证；
  \item Sign$_{sk}(x)$：私钥签名、公钥验证，支持第三方验证；“不可否认性”还依赖密钥控制、审计和法律/制度语境，不能由签名算法单独保证；
  \item 三者都不等同于“把内容隐藏起来”。
\end{itemize}
\vfill
\begin{center}
\begin{tikzpicture}[>=Latex,node distance=10mm]
\nodebox[minimum width=2.3cm]{ADSBlue}{msg}{消息};
\nodebox[right=of msg,minimum width=2.3cm]{ADSOrange}{hash}{摘要/MAC/签名};
\nodebox[right=of hash,minimum width=2.3cm]{ADSGreen}{ver}{验证结果};
\draw[\arrow,ADSGray] (msg)--(hash)--(ver);
\end{tikzpicture}
\end{center}
\end{frame}

\begin{frame}{认证协议的三条硬约束}
\begin{enumerate}
  \item \key{来源}：对方确实掌握与身份绑定的秘密；
  \item \key{新鲜性}：消息是这次会话产生的，不是旧包重放；
  \item \key{绑定}：挑战、接收者、用途和上下文都被纳入保护范围。
\end{enumerate}
\vfill
\callout{反例}{只加 nonce 但不把“接收者 Q”放入 MAC，可能被中间人转发给另一个服务。}
\end{frame}

\sectionframe{三、共享密钥认证：从脆弱到可用}{先看最简单的协议，再一步步指出它为什么不够。}

\begin{frame}{基本共享密钥协议}
设 P、Q 共享密钥 $k$：
\begin{enumerate}
  \item P 发送明文 $m=\text{“I am P”}$；
  \item P 发送 $m' = \{m,Q\}_k$；
  \item Q 用 $k$ 重算并比较 $m'$。
\end{enumerate}
\begin{center}
\begin{tikzpicture}[>=Latex,node distance=12mm]
\nodebox[minimum width=2.3cm]{ADSBlue}{p}{P};
\nodebox[right=of p,minimum width=2.3cm]{ADSOrange}{q}{Q};
\draw[\arrow,ADSGray] (p)--node[above,font=\scriptsize]{m,m'}(q);
\draw[ADSGray,dashed] ($(p.south)!0.5!(q.south)+(0,-8mm)$) node[font=\scriptsize]{Q 验证 $\{m,Q\}_k=m'$};
\end{tikzpicture}
\end{center}
\end{frame}

\begin{frame}{基本协议为什么会被重放？}
\begin{itemize}
  \item 攻击者记录一次合法的 $(m,m')$；
  \item 未来再次发送给 Q；
  \item Q 无法区分“刚刚生成的证明”和“历史录像”；
  \item 若消息触发转账、删除或调度，旧证明就可能造成真实损害。
\end{itemize}
\vfill
\callout{关键观察}{认证不仅是“证明知道秘密”，还要证明“现在正在使用秘密”。}
\end{frame}

\begin{frame}{重放攻击完整例子：旧证明为什么仍会通过}
\small
设 $m=(P,Q,\mathrm{transfer}(100))$，认证标签为 $t=\mathrm{MAC}_k(m)$。
\begin{center}
\begin{tabularx}{.96\linewidth}{>{\bfseries}p{1.5cm}p{3.5cm}X}
\toprule
阶段 & 消息 & Q 的判断 \\
\midrule
合法会话 & $P\rightarrow Q:(m,t)$ & MAC 验证通过，Q 执行一次转账。 \\
窃听记录 & 攻击者 E 保存 $(m,t)$ & E 不需要知道密钥，也不需要修改消息。 \\
稍后重放 & $E\rightarrow Q:(m,t)$ & 同一 MAC 仍然正确；协议中没有会话号、nonce 或时间边界。 \\
错误结果 & Q 再次接受 & Q 无法知道这是“历史录像”，动作可能被重复执行。 \\
\bottomrule
\end{tabularx}
\end{center}
\callout{攻击成立的原因}{MAC 证明消息来自密钥持有者且未被篡改，但它本身不证明消息是本轮新生成的。}
\end{frame}

\begin{frame}{Nonce 挑战—响应协议}
\begin{enumerate}
  \item P $\rightarrow$ Q：我是 P；
  \item Q 生成大空间随机数 $n$，发送给 P；
  \item P $\rightarrow$ Q：$\{P,Q,n\}_k$；
  \item Q 验证密文，接受本轮身份。
\end{enumerate}
\begin{center}
\begin{tikzpicture}[>=Latex,node distance=10mm]
\nodebox[minimum width=2.3cm]{ADSBlue}{p2}{P};
\nodebox[right=of p2,minimum width=2.3cm]{ADSOrange}{q2}{Q};
\draw[\arrow,ADSGray] (p2)--node[above,font=\scriptsize]{P}(q2);
\draw[\arrow,ADSBlue] (q2)--node[below,font=\scriptsize]{nonce $n$}(p2);
\draw[\arrow,ADSGreen] (p2)--node[above,font=\scriptsize]{$\{P,Q,n\}_k$}(q2);
\end{tikzpicture}
\end{center}
\end{frame}

\begin{frame}{Nonce 修复完整例子：旧响应如何被拒绝}
\small
令 $a=\mathrm{transfer}(100)$。把身份、接收者、动作和本轮挑战一起绑定：
\[
t_1=\mathrm{MAC}_k(P,Q,n_1,a).
\]
\begin{center}
\begin{tabularx}{.96\linewidth}{>{\bfseries}p{1.3cm}>{\raggedright\arraybackslash}p{4.2cm}>{\raggedright\arraybackslash}X}
\toprule
步骤 & 消息/状态 & 验证结果 \\
\midrule
1 & Q 生成未使用的随机数 $n_1$ & Q 把 $n_1$ 标记为当前待完成挑战。 \\
2 & $Q\rightarrow P:n_1$ & P 为本轮挑战计算响应 $t_1$。 \\
3 & $P\rightarrow Q:(P,n_1,a,t_1)$ & Q 验证 MAC 与待完成的 $n_1$，接受后立即把 $n_1$ 标记为已使用。 \\
4 & E 重放旧的 $(n_0,a,t_0)$ & $n_0$ 不是当前挑战，拒绝。 \\
5 & E 再次重放 $(n_1,a,t_1)$ & $n_1$ 已消费，拒绝。 \\
\bottomrule
\end{tabularx}
\end{center}
\smallnote{Nonce 解决认证新鲜性；业务重试仍需 request ID、幂等写或事务机制保证“只执行一次”。}
\end{frame}

\begin{frame}{Nonce 协议的边界}
\begin{itemize}
  \item nonce 必须不可预测、足够长且每轮只用一次；
  \item Q 必须把 Q 自己也放入保护内容，防止跨服务转发；
  \item P、Q 需要知道对方的长期密钥，规模为 $O(n^2)$ 时难以维护；
  \item 仍需考虑密钥轮换、失效和审计。
\end{itemize}
\end{frame}

\begin{frame}{时间戳与 nonce：两种新鲜性来源}
\begin{tabularx}{\linewidth}{>{\bfseries}p{3cm}XX}
\toprule
Nonce & 不依赖全局时钟；需要保存/检测挑战；适合异步网络。 & 服务器有状态，需防止 nonce 表膨胀。 \\
时间戳 & 服务器可快速判断过期；票据可设置生命周期。 & 依赖时钟同步；时钟漂移会误拒绝。 \\
混合 & 时间窗口 + nonce + 会话标识。 & 协议更健壮但实现更复杂。 \\
\bottomrule
\end{tabularx}
\vfill
\callout{课堂判断}{若服务跨越多个时区且网络延迟长，优先选择哪一种？先说出隐含假设，再做选择。}
\end{frame}

\begin{frame}{Wide-Mouth Frog：让认证服务器转发会话密钥}
\begin{itemize}
  \item A 与认证服务器 S 共享 $K_{AS}$，B 与 S 共享 $K_{BS}$；
  \item A 生成会话密钥 $K_{AB}$，请求 S 转交给 B；
  \item S 发送包含时间戳、A 身份和 $K_{AB}$ 的密文给 B；
  \item 优点：A 不必与 B 预先共享长期密钥；缺点：依赖时钟，S 被攻破会影响全局。
\end{itemize}
\end{frame}

\begin{frame}{认证服务器协议：可扩展，但中心更重要}
\begin{center}
\resizebox{.95\linewidth}{!}{%
\begin{tikzpicture}[>=Latex,node distance=16mm]
\nodebox[minimum width=2.1cm]{ADSBlue}{pa}{P};
\nodebox[right=of pa,minimum width=2.3cm]{ADSOrange}{qa}{Q};
\nodebox[below=of pa,minimum width=2.4cm]{ADSGreen}{sa}{认证服务器 S};
\draw[\arrow,ADSGray] (pa)--node[above,font=\scriptsize]{身份}(qa);
\draw[\arrow,ADSBlue] (qa)--node[right,font=\scriptsize]{挑战 n}(pa);
\draw[\arrow,ADSGreen] (pa)--node[left,font=\scriptsize]{$\{P,Q,n\}_{K_{PS}}$}(sa);
\draw[\arrow,ADSGreen] (qa)--node[right,font=\scriptsize]{$\{P,Q,x\}_{K_{QS}}$}(sa);
\draw[\arrow,ADSOrange] (sa)--node[below,font=\scriptsize]{给 Q 返回可验证结果}(qa);
\end{tikzpicture}}
\end{center}
\vspace{-6mm}
\smallnote{每个 principal 只需与 S 共享一个密钥，但 S 成为可用性与信任根。}
\end{frame}

\sectionframe{四、无时间戳的协议与 Kerberos}{认证服务器不只是“查密码”，还负责把信任转换成短期票据。}

\begin{frame}{Otway--Rees 协议的角色}
\begin{itemize}
  \item A、B 各自与 S 共享 $K_{AS}$、$K_{BS}$；
  \item $N_A,N_B$ 是双方 nonce，M 是事务标识；
  \item S 生成会话密钥 $K_{AB}$，分别加密给 A、B；
  \item 不需要全局时间戳，但必须严谨绑定事务和双方身份。
\end{itemize}
\end{frame}

\begin{frame}{Otway--Rees 四步消息}
\begin{enumerate}
  \item A $\rightarrow$ B：$M,A,B,\{N_A,M,A,B\}_{K_{AS}}$；
  \item B $\rightarrow$ S：追加 $\{N_B,M,A,B\}_{K_{BS}}$；
  \item S $\rightarrow$ B：分别返回 $\{N_A,K_{AB}\}_{K_{AS}}$ 与 $\{N_B,K_{AB}\}_{K_{BS}}$；
  \item B $\rightarrow$ A：转发 A 的密文部分。
\end{enumerate}
\vfill
\callout{分析任务}{第 4 步被拦截并重放时，A、B 是否一定得到同一个会话密钥？为什么？}
\end{frame}

\begin{frame}{Kerberos：把长期密码换成短期票据}
\begin{itemize}
  \item 用户只在登录阶段证明密码；
  \item 认证服务器签发 TGT（ticket-granting ticket）；
  \item TGS 根据 TGT 签发目标服务票据；
  \item 客户端携带服务票据访问服务，服务端无需知道用户密码。
\end{itemize}
\begin{center}
\begin{tikzpicture}[>=Latex,node distance=10mm and 14mm]
\nodebox[minimum width=2cm]{ADSBlue}{c}{客户端};
\nodebox[right=of c,minimum width=2cm]{ADSOrange}{as}{AS};
\nodebox[right=of as,minimum width=2cm]{ADSGreen}{tgs}{TGS};
\nodebox[right=of tgs,minimum width=2cm]{ADSGray}{ss}{服务 S};
\draw[\arrow,ADSGray] (c)--node[above,font=\scriptsize]{登录}(as);
\draw[\arrow,ADSGray] (c)--node[below,font=\scriptsize]{TGT}(tgs);
\draw[\arrow,ADSGray] (c)--node[above,font=\scriptsize]{服务票据}(ss);
\end{tikzpicture}
\end{center}
\end{frame}

\begin{frame}{Kerberos 初始登录}
\begin{enumerate}
  \item U 向主机 H 声明身份；
  \item H 向 Kerberos 请求 U 的登录材料；
  \item Kerberos 生成会话密钥 $k$ 和 TGT；
  \item 结果用由用户密码导出的 $k_U$ 加密返回 H；
  \item 用户输入密码，H 解密并立即擦除密码。
\end{enumerate}
\vfill
\smallnote{密码不应被长期存储在应用内存中；长期信任由 KDC 和票据生命周期承担。}
\end{frame}

\begin{frame}{Kerberos 获取服务票据}
\begin{itemize}
  \item 客户端把 TGT、目标服务 S 和时间戳 $T_1$ 发给 TGS；
  \item TGS 用自己的密钥解开 TGT，验证时间窗口；
  \item TGS 生成新的客户端—服务会话密钥 $k'$；
  \item 返回用旧会话密钥加密的结果和服务票据 $\{C,S,k',T,L\}_{K_S}$。
\end{itemize}
\end{frame}

\begin{frame}{Kerberos 请求服务与双向认证}
\begin{enumerate}
  \item C $\rightarrow$ S：服务票据 $ticket_S$ 与认证器 $\{C,T_2\}_{k'}$；
  \item S 用 $K_S$ 解票据，得到 $k'$；
  \item S 用 $k'$ 解认证器，检查 C 与时间戳；
  \item 若需要双向认证，S 返回 $\{T_2+1\}_{k'}$。
\end{enumerate}
\vfill
\callout{工程含义}{服务端只需保护自己的长期密钥，用户密码不在服务间流动。}
\end{frame}

\begin{frame}{Kerberos 的现实边界}
\begin{itemize}
  \item KDC 是信任根：必须高可用、可复制、可审计；
  \item 时间戳要求时钟同步，时钟漂移会造成误判；
  \item 票据泄露后，在有效期内可能被冒用；
  \item 跨域、跨云和离线设备需要额外的信任桥接。
\end{itemize}
\end{frame}

\sectionframe{五、一次性口令、证书与现代身份}{把“拥有秘密”升级为“可撤销、可轮换、可审计的凭证”。}

\begin{frame}{一次性口令（OTP）的核心}
\begin{itemize}
  \item 注册阶段，客户端与服务器共享种子或口令链；
  \item 第 $t$ 次登录只使用第 $t$ 个口令；
  \item 服务器保存链上的下一个验证值，而不是明文口令；
  \item 即使历史口令泄露，重放也不能通过下一轮验证。
\end{itemize}
\end{frame}

\begin{frame}{Hash chain：为什么服务器可以只存一个值？}
设 $p_i=H(p_{i+1})$：
\begin{itemize}
  \item 服务器存 $p_0$；客户端第一次提交 $p_1$；
  \item 服务器计算 $H(p_1)$ 并与 $p_0$ 比较；
  \item 验证成功后把 $p_1$ 作为新的基准；
  \item 单向性使服务器难以从旧值反推出未来口令。
\end{itemize}
\vfill
\callout{注意}{OTP 主要解决重放，不自动解决钓鱼；攻击者仍可能实时转发一次性口令。}
\end{frame}

\begin{frame}{公开密钥与证书链}
\begin{center}
\begin{tikzpicture}[>=Latex,node distance=11mm and 14mm]
\nodebox[minimum width=2.2cm]{ADSBlue}{root}{根 CA};
\nodebox[below left=of root,minimum width=2.5cm]{ADSOrange}{ica}{中间 CA};
\nodebox[below right=of root,minimum width=2.5cm]{ADSGreen}{svc1}{服务证书};
\draw[\arrow,ADSGray] (root)--(ica);
\draw[\arrow,ADSGray] (ica)--(svc1);
\node[below=of svc1,font=\scriptsize,align=center]{公钥 + 主体 + 有效期\\+ 用途 + CA 签名};
\end{tikzpicture}
\end{center}
\begin{itemize}
  \item 验证者用上级 CA 公钥验证下级证书；
  \item 证书把“公钥”绑定到 DNS 名称、服务身份或组织；
  \item 撤销列表、短有效期和自动续期决定运营成本。
\end{itemize}
\end{frame}

\begin{frame}{签名挑战—响应}
\begin{enumerate}
  \item Q 发送随机挑战 $n$ 和上下文 $ctx$；
  \item P 返回 $sig=Sign_{sk_P}(P,Q,n,ctx)$；
  \item Q 用证书中的 $pk_P$ 验签并检查证书状态；
  \item 验证成功后把身份映射到授权策略。
\end{enumerate}
\vfill
\smallnote{把 Q、用途和会话标识放入签名，防止“同一签名被拿去别的服务使用”。}
\end{frame}

\begin{frame}{mTLS：服务网格里的双向认证}
\begin{itemize}
  \item 客户端与服务端都出示证书；
  \item 握手阶段验证证书链、域名/身份和有效期；
  \item 双方协商短期会话密钥，后续请求用对称加密保护；
  \item 代理可以把证书主体映射成 service account 和 RBAC 规则。
\end{itemize}
\end{frame}

\begin{frame}{近期案例：SPIFFE/SPIRE 工作负载身份}
\begin{itemize}
  \item 工作负载启动后由身份系统依据命名、节点和证明材料签发短期 SVID，而不是把长期私钥写进镜像；
  \item 服务间使用 mTLS 验证工作负载身份，授权策略再根据身份、资源和动作做决定；
  \item 证书自动轮换、过期和撤销把“凭证生命周期”纳入控制环；
  \item 对 AI 平台，训练作业、模型仓、数据仓和 GPU 控制器可以各自拥有独立 workload identity，降低共享管理员密钥的横向风险。
\end{itemize}
\vfill
\smallnote{本页关注身份生命周期的分布式系统抽象；具体部署仍需明确 attestation、信任根、轮换窗口和离线恢复策略。}
\end{frame}

\sectionframe{六、认证之后：授权与零信任}{“身份正确”只是入口；每个请求仍要依据最小权限和上下文作决定。}

\begin{frame}{RBAC、ABAC 与能力令牌}
\begin{tabularx}{\linewidth}{>{\bfseries}p{3cm}XX}
\toprule
RBAC & 用户/服务绑定角色，角色绑定权限。 & 易审计，角色爆炸时管理困难。 \\
ABAC & 根据主体、资源、动作、环境属性决策。 & 灵活，策略解释与测试更复杂。 \\
能力令牌 & 持有不可伪造令牌即拥有有限能力。 & 便于委托，但泄露后需快速撤销。 \\
\bottomrule
\end{tabularx}
\vfill
\callout{例子}{“训练作业只能读取数据集版本 v7、写入实验桶，不能访问生产模型”是资源级授权，而不是简单角色名。}
\end{frame}

\begin{frame}{零信任的请求路径}
\begin{center}
\resizebox{.96\linewidth}{!}{%
\begin{tikzpicture}[>=Latex,node distance=12mm and 12mm]
\nodebox[minimum width=2.1cm]{ADSBlue}{clientz}{调用方};
\nodebox[right=of clientz,minimum width=2.4cm]{ADSOrange}{proxyz}{策略代理};
\nodebox[right=of proxyz,minimum width=2.2cm]{ADSGreen}{pdp}{策略决策};
\nodebox[right=of pdp,minimum width=2.2cm]{ADSGray}{resource}{资源服务};
\draw[\arrow,ADSGray] (clientz)--node[above,font=\scriptsize]{身份+请求}(proxyz);
\draw[\arrow,ADSBlue] (proxyz)--node[above,font=\scriptsize]{属性/策略查询}(pdp);
\draw[\arrow,ADSGreen] (pdp)--node[above,font=\scriptsize]{allow/deny}(proxyz);
\draw[\arrow,ADSOrange] (proxyz)--node[above,font=\scriptsize]{已授权请求}(resource);
\end{tikzpicture}}
\end{center}
\begin{itemize}
  \item 不默认信任网络位置；每次请求重新验证身份和策略；
  \item 审计日志必须记录主体、资源、动作、结果和策略版本。
\end{itemize}
\end{frame}

\begin{frame}{密钥轮换与撤销}
\begin{enumerate}
  \item 生成新密钥并标记 key-id；
  \item 双读旧/新、单写新，平滑迁移存量数据；
  \item 过渡窗口结束后撤销旧 key-id；
  \item 监测异常使用，必要时立即吊销并回滚依赖。
\end{enumerate}
\vfill
\callout{常见事故}{把密钥轮换理解为“替换一个文件”，忽略了缓存、长连接、备份和异步任务，最终造成半数请求失败。}
\end{frame}

\sectionframe{七、AI 基础设施中的身份与安全}{模型、数据、GPU 和控制面都需要可验证的主体与最小权限。}

\begin{frame}{AI 训练平台的主体关系}
\begin{center}
\begin{tikzpicture}[>=Latex,node distance=10mm and 13mm]
\nodebox[minimum width=2.2cm]{ADSBlue}{job}{训练作业};
\nodebox[right=of job,minimum width=2.2cm]{ADSOrange}{sched}{调度器};
\nodebox[below=of job,minimum width=2.2cm]{ADSGreen}{data}{数据仓};
\nodebox[below=of sched,minimum width=2.2cm]{ADSGray}{model}{模型仓};
\draw[\arrow,ADSGray] (job)--node[above,font=\scriptsize]{申请 GPU}(sched);
\draw[\arrow,ADSBlue] (job)--node[left,font=\scriptsize]{读训练数据}(data);
\draw[\arrow,ADSOrange] (job)--node[right,font=\scriptsize]{写 checkpoint}(model);
\end{tikzpicture}
\end{center}
\begin{itemize}
  \item 每个作业使用独立 workload identity，而不是共享管理员密钥；
  \item 数据、模型、镜像、GPU 资源分别定义权限边界。
\end{itemize}
\end{frame}

\begin{frame}{模型仓库：身份、完整性和版本}
\begin{itemize}
  \item 上传者签名模型清单、权重摘要和训练来源；
  \item 读取者验证签名、版本和审批状态；
  \item 生产部署只允许来自受信任 registry 的不可变 digest；
  \item 回滚时保留旧版本的审计记录，而不是覆盖标签。
\end{itemize}
\vfill
\callout{安全属性}{“能下载模型”不等于“模型来自可信来源”，因此需要签名、SBOM、扫描结果和发布策略。}
\end{frame}

\begin{frame}{多租户 GPU 集群：隔离不止是命名空间}
\begin{itemize}
  \item 控制面：谁能创建作业、抢占 GPU、读取日志；
  \item 数据面：GPU 显存、共享文件、缓存和 NCCL 通道是否跨租户泄露；
  \item 网络面：训练作业能否访问 metadata server、KDC 或其他租户服务；
  \item 供应链：基础镜像、驱动、插件和模型代码是否可验证。
\end{itemize}
\end{frame}

\begin{frame}{数据和权重的保护路径}
\begin{center}
\begin{tikzpicture}[>=Latex,node distance=12mm and 14mm]
\nodebox[minimum width=2.0cm]{ADSBlue}{ing}{数据入口};
\nodebox[right=of ing,minimum width=2.2cm]{ADSOrange}{enc}{加密存储};
\nodebox[right=of enc,minimum width=2.2cm]{ADSGreen}{train}{训练作业};
\nodebox[right=of train,minimum width=2.2cm]{ADSGray}{reg}{模型仓};
\draw[\arrow,ADSGray] (ing)--node[above,font=\scriptsize]{校验/脱敏}(enc);
\draw[\arrow,ADSBlue] (enc)--node[above,font=\scriptsize]{短期凭证}(train);
\draw[\arrow,ADSOrange] (train)--node[above,font=\scriptsize]{签名权重}(reg);
\end{tikzpicture}
\end{center}
\begin{itemize}
  \item KMS 只向有权限的作业发放短期解密能力；
  \item 日志中记录 key-id、模型 digest 和作业身份，便于追溯。
\end{itemize}
\end{frame}

\begin{frame}{AI 服务的请求授权}
\begin{itemize}
  \item 客户端证书确认服务调用方；
  \item token 中绑定模型名、租户、配额、请求范围和过期时间；
  \item 网关检查速率、预算、地域和敏感数据策略；
  \item 推理服务只接收已验证上下文，不自行信任用户传入的 tenant-id。
\end{itemize}
\end{frame}

\begin{frame}{供应链签名：为什么要签“整个构建结果”？}
\begin{itemize}
  \item 只签源码，无法证明实际运行的镜像、依赖和编译器；
  \item 对镜像 digest、模型权重、配置和 SBOM 做统一签名；
  \item 部署准入控制器验证签名者、构建流水线和漏洞门槛；
  \item 运行时仍需监控，因为合法构建也可能被错误配置滥用。
\end{itemize}
\end{frame}

\sectionframe{八、课堂推演与协议审查}{用攻击者视角检查“协议写出来了”是否等于“协议安全”。}

\begin{frame}{互动一：找出重放窗口}
\begin{block}{协议}
P $\rightarrow$ Q：$P,\{P,\text{transfer}=100\}_{k}$；\quad Q：验证后执行转账。
\end{block}
\begin{enumerate}
  \item 攻击者记录合法消息；
  \item 在一小时后重复发送；
  \item 小组回答：nonce、时间戳、序列号分别如何修改协议？
\end{enumerate}
\vfill
\smallnote{建议答案：将接收者、操作、金额和一次性挑战绑定在 MAC 中；服务端拒绝已消费挑战。}
\end{frame}

\begin{frame}{互动二：Kerberos 票据失窃}
\begin{itemize}
  \item 假设攻击者窃取了有效期 8 小时的服务票据；
  \item 服务端没有在线查询撤销状态；
  \item 讨论：缩短票据、绑定终端、增加二次证明、快速轮换，哪一种最有效？
\end{itemize}
\vfill
\callout{评分点}{回答必须同时谈到安全性、可用性和认证服务器压力。}
\end{frame}

\begin{frame}{互动三：给训练作业设计最小权限}
\begin{tabularx}{\linewidth}{>{\bfseries}p{3.1cm}XX}
\toprule
主体 & 允许 & 禁止 \\
数据清洗作业 & 读取原始数据、写入脱敏桶 & 读取模型仓生产版本 \\
训练作业 & 读取脱敏桶、写入实验模型 & 删除数据、访问其他租户日志 \\
发布流水线 & 读取已审批权重、写入生产 registry & 直接修改训练数据 \\
\bottomrule
\end{tabularx}
\vfill
\smallnote{先列资源和动作，再选择 RBAC/ABAC；不要从“给某个团队管理员权限”开始。}
\end{frame}

\begin{frame}{协议审查清单}
\begin{enumerate}
  \item 谁是 principal？密钥/证书如何绑定到它？
  \item 消息是否包含接收者、用途、会话和版本？
  \item 新鲜性由 nonce、时间戳还是序列号提供？
  \item 失败、超时、重试、并发登录时会发生什么？
  \item 凭证如何轮换、撤销、审计和恢复？
\end{enumerate}
\end{frame}

\sectionframe{九、总结}{认证是分布式系统的信任根，安全性来自协议细节和运维闭环。}

\begin{frame}{本讲知识地图}
\begin{center}
\begin{tikzpicture}[>=Latex,node distance=11mm and 14mm]
\nodebox[minimum width=2.2cm]{ADSBlue}{goal}{安全目标};
\nodebox[below left=of goal,minimum width=2.3cm]{ADSOrange}{sym}{共享密钥/nonce};
\nodebox[below right=of goal,minimum width=2.3cm]{ADSGreen}{pub}{公钥/证书};
\nodebox[below=of sym,minimum width=2.5cm]{ADSGray}{ticket}{票据/Kerberos};
\nodebox[below=of pub,minimum width=2.5cm]{ADSGray}{policy}{授权/零信任};
\nodebox[below=of ticket,minimum width=2.6cm]{ADSBlue}{ai}{AI 基础设施};
\draw[\arrow,ADSGray] (goal)--(sym);
\draw[\arrow,ADSGray] (goal)--(pub);
\draw[\arrow,ADSGray] (sym)--(ticket);
\draw[\arrow,ADSGray] (pub)--(policy);
\draw[\arrow,ADSGray] (ticket)--(ai);
\draw[\arrow,ADSGray] (policy)--(ai);
\end{tikzpicture}
\end{center}
\vspace{-5mm}
\smallnote{带走一句话：先认证“谁”，再授权“做什么”；证明必须绑定新鲜性、上下文和生命周期。}
\end{frame}

\begin{frame}{课后练习与阅读}
\begin{enumerate}
  \item 把基本共享密钥协议改写成支持双向认证的消息序列；
  \item 画出 Kerberos 票据从登录到服务访问的状态转换；
  \item 为一个多租户推理服务写出 5 条最小权限策略；
  \item 比较 Otway--Rees、认证服务器和 Kerberos 的假设差异，并画出消息时序。
\end{enumerate}
\vfill
\smallnote{下一讲：P2P、DHT、自稳定与大规模 AI 基础设施的系统综合。}
\end{frame}

\begin{frame}{安全工具的边界：证明身份不等于获得权限}
\small
\begin{itemize}
  \item Hash 只能提供摘要；MAC 证明持有共享密钥的一方生成了内容；数字签名支持公开验证，但这些都不自动给出业务授权。
  \item 证书证明“公钥与某个身份绑定”，不等于该身份可以访问任意租户、模型或数据；授权还要结合策略、资源和操作上下文。
  \item nonce 防止旧消息直接重放，但还必须绑定接收者、动作、租户、版本和有效期，否则可能发生跨服务转发或语义重放。
  \item Kerberos、mTLS 和服务账号都依赖密钥轮换、撤销、时钟/票据窗口和最小权限；协议图省略的生命周期管理，往往决定真实安全性。
\end{itemize}
\vfill
\callout{安全审计顺序}{先问“谁签了”，再问“签了什么”，最后问“策略是否允许这次操作”；三者不能用一个认证结果代替。}
\end{frame}

\end{document}
